\begin{table}[htbp]
\centering
\caption{Border premium, free entry (New Mexico) vs.\ capped entry (Colorado + Montana), 100 miles}
\label{tab:county_border_capped_nmcomt_100mi}
\footnotesize
\begin{tabular}{lccccc}
\toprule
& \multicolumn{2}{c}{Border (NM, free)} & \multicolumn{2}{c}{Border $\times$ Capped (CO+MT vs.\ NM)} & \\
Outcome & Coef. (\%) & $p$ (cl.) & Coef. (\%) & $p$ (cl.) & $n$ \\
\midrule
Revenue per capita & $+0.684$ ($+98.2\%$) & 0.001 & $-0.507$ ($-39.8\%$) & 0.069 & 2478 \\
Revenue per store & $+0.088$ ($+9.2\%$) & 0.655 & $-0.249$ ($-22.1\%$) & 0.296 & 1007 \\
Dispensary density & $+0.596$ ($+81.5\%$) & 0.012 & $-0.357$ ($-30.0\%$) & 0.202 & 1007 \\
\bottomrule
\end{tabular}
\begin{minipage}{0.98\textwidth}
\vspace{0.3em}
\footnotesize
\textit{Notes:} Sample restricted to New Mexico (always free entry) and Colorado + Montana (always capped
entry) -- Oregon excluded entirely, since it switches regime mid-panel and would contribute to both sides of
this comparison. Border = county centroid within 100 miles of the nonlegal-state line, with
$\geq$1,000 nonlegal-state residents within that reach. State and state $\times$ quarter fixed effects via a
single concatenated state-quarter key. County-clustered standard errors.
\end{minipage}
\end{table}
